% 小熊座（综述）
% license CCBYSA3
% type Wiki

本文根据 CC-BY-SA 协议转载翻译自维基百科\href{https://en.wikipedia.org/wiki/Ursa_Minor}{相关文章}。

小熊座（Ursa Minor，拉丁语意为“较小的熊”，与大熊座 Ursa Major 相对），又称小熊星座，是位于遥远北天的一组星座。与大熊座类似，小熊座的尾部也可被视作一只勺子的把手，因此在北美被称为“小北斗”（Little Dipper）：由七颗恒星组成，其中四颗构成勺状的斗部，形态与其对应的大北斗相似。小熊座是公元$^{[2]}$世纪天文学家托勒密列出的$^{[48]}$个星座之一，至今仍是现代$^{[88]}$个星座之一。由于北极星位于其中，小熊座在传统航海导航中，尤其对水手而言，具有重要意义。

北极星（Polaris）是该星座中最亮的恒星，为一颗黄白色超巨星，同时也是夜空中最亮的造父变星，其视星等在$^{[1.97]}$至$^{[2.00]}$之间变化。小熊座$\beta$星（Beta Ursae Minoris），又名 Kochab，是一颗已进入晚期演化阶段的恒星，因膨胀并冷却而成为一颗橙色巨星，视星等为$^{[2.08]}$，仅略暗于北极星。Kochab 与视星等约为$^{[3]}$等的小熊座$\gamma$星（Gamma Ursae Minoris）常被称为“北极星的守护者”或“极星守卫者”$^{[4]}$。在该星座中已有行星被探测到环绕其中四颗恒星运行，包括 Kochab。此外，小熊座还包含一颗孤立的中子星——Calvera，以及迄今发现的第二高温白矮星 H1504+65，其表面温度约为$^{[200{,}000]}$ K。
\subsection{历史与神话}
\begin{figure}[ht]
\centering
\includegraphics[width=6cm]{./figures/9a220c9be61db089.png}
\caption{天龙座环绕着小熊座，如约 1825 年在伦敦出版的一组星座图《乌拉尼娅之镜》（Urania's Mirror）中所描绘的那样。} \label{fig_Ursa_1}
\end{figure}
在巴比伦的星表中，小熊座被称为“天之车”（MUL\,MAR.GÍD.DA.AN.NA，并与女神 Damkina 相关联）。它被列入约公元前 1000 年编纂的《MUL.APIN》星表之中，属于“恩利尔之星”，即北天星区。$^{[6]}$ 据第欧根尼·拉尔修记载并引述卡利马科斯的话，米利都的泰勒斯“测量了腓尼基人航海所依凭的那辆车的星辰”。第欧根尼将其认定为小熊座；由于据说腓尼基人在海上航行时使用该星座，它也被称为 Phoinikē。$^{[7][8]}$ 将北天星座称作“熊”的传统似乎确实源自希腊，尽管荷马只提到过一只“熊”。$^{[9]}$ 因而最初的“熊”应为大熊座；而小熊座作为第二只熊，或“腓尼基之熊”（Ursa Phoenicia，即 Φοινίκη），则是后来才被接纳的。根据斯特拉波（I.1.6, C3）的说法，这是由于泰勒斯的建议，他向希腊人推荐该星座作为航海指引，而此前希腊人一直以大熊座进行航行定位。在古典时代，天极的位置略更接近小熊座的 $\beta$ 星而非 $\alpha$ 星，因此整个星座都被视为指示北方。自中世纪起，使用小熊座 $\alpha$ 星（即“北极星”）作为北极指向变得更加便利，尽管在中世纪北极星仍与天极相距数度。$^{[10][a]}$ 如今，北极星已位于距离北天极 1° 以内，并仍然是当前的极星。其新拉丁名称 \emph{stella polaris} 仅在近代早期才出现。$^{[4]}$  

该星座的另一古名为 Cynosura（希腊语 Κυνοσούρα，“狗之尾”）。这一名称的起源并不明确（若将小熊座视为“狗尾”，则意味着附近应存在一只“狗”的星座，但并无此星座）。$^{[11]}$ 相反，在《星变志》（Catasterismi）的神话传统中，Cynosura 被解释为一位俄瑞阿得斯山林仙女，她曾担任宙斯的乳母，因而受到宙斯的褒奖，被安置于天穹之中。$^{[12]}$ 关于 Cynosura 这一名称还存在多种解释：有一种将其与卡利斯托的神话联系起来，认为其子阿耳卡斯被替换为她的狗，并由宙斯将其置于天上；$^{[11]}$ 也有人认为在更早的观念中，大熊座被视为一头牛，与牧人座构成牧牛组合，而小熊座则是一只狗。$^{[13]}$ 乔治·威廉·考克斯将该名称解释为 Λυκόσουρα 的变体，理解为“狼之尾”，但他本人又将其词源解释为“光的轨迹或行列”（即 λύκος“狼”与 λύκ-“光”之间的关联）。艾伦则指出，可将其与古爱尔兰语中该星座的名称 drag-blod（“火之轨迹”）进行比较。布朗（1899）提出 Cynosura 可能并非希腊语起源，而是借自亚述语 An-nas-sur-ra，意为“高升者”。$^{[14]}$  

另有一种神话叙述两只熊曾在伊达山上将宙斯隐藏起来，使其免遭其父克洛诺斯的杀害；后来宙斯将它们置于天空之中，但由于被神抛向天穹，它们的尾巴被拉得很长。$^{[15]}$  

由于小熊座由七颗恒星组成，拉丁语中表示“北方”的词 \emph{septentrio}（即北极星所指的方向）源自 \emph{septem}（七）与 \emph{triones}（牛），意为“拉犁的七头牛”，而这七颗星也常被想象为这种形象。该名称同样被用来指称大熊座的主星。$^{[16]}$  

在因纽特天文学中，最亮的三颗星——北极星、Kochab 与 Pherkad——被称为 Nuutuittut（“永不移动者”），尽管这一称呼更常以单数形式使用，专指北极星。由于在极北纬地区北极星在天空中的高度过高，它并不适合用于航海定位。$^{[17]}$ 在中国天文学中，小熊座的主要恒星被划分为两个星官：勾陈（Gòuchén，包含 $\alpha$、$\delta$、$\varepsilon$、$\zeta$、$\eta$、$\theta$、$\lambda$ UMi）以及北极（Běijí，包含 $\beta$ 与 $\gamma$ UMi）。$^{[18]}$  
\subsection{特征}
\begin{figure}[ht]
\centering
\includegraphics[width=8cm]{./figures/a72e9261fedb4cea.png}
\caption{约公元 $^{[1009-1010]}$ 年《定星书》中所描绘的小熊座。与西方的表现方式不同，图中的熊被画成尾巴下垂的形态。} \label{fig_Ursa_2}
\end{figure}
小熊座西侧与鹿豹座相邻，西侧亦与天龙座接壤，东侧则与仙王座相邻。其覆盖面积为 $256$ 平方度，在 $88$ 个星座中按面积大小排名第 $56$ 位。小熊座在美国的通俗称呼为“小北斗”，因为其最亮的七颗恒星似乎构成了一个长柄勺（舀勺）的形状。勺柄末端的恒星是北极星 Polaris。北极星也可以通过沿着大熊座“斗碗”末端形成的两颗恒星——大熊座 $\alpha$ 星与 $\beta$ 星（通常称为“指极星”）连成的直线，在夜空中延伸约 $30^\circ$（相当于手臂伸直时三个竖起的拳头的宽度）来找到 $^{[19]}$。构成小北斗“斗碗”的四颗恒星，其视星等分别为二等、三等、四等和五等星，为判断肉眼可见恒星的星等提供了一个简便的参照，对城市居民或用于测试视力尤为有用 $^{[20]}$。

该星座的三字母缩写由国际天文学联合会（IAU）于 $^{[1922]}$ 年正式采用，为“UMi”$^{[21]}$。官方星座边界由比利时天文学家欧仁·德尔波特于 $^{[1930]}$ 年确定，并以一个由 $22$ 条边组成的多边形来界定（见信息框示意）。在赤道坐标系中，这些边界的赤经范围介于 $08^\mathrm{h}\,41.4^\mathrm{m}$ 与 $22^\mathrm{h}\,54.0^\mathrm{m}$ 之间，而赤纬范围则从北天极延伸至南侧的 $65.40^\circ$ $^{[1]}$。由于其位于遥远的北天半球，小熊座整体仅能被北半球的观测者完整看到 $^{[22][b]}$。
\subsection{特征}
\begin{figure}[ht]
\centering
\includegraphics[width=6cm]{./figures/9675489a80e4d9ed.png}
\caption{肉眼所见的小熊座星座（已添加连线与标注）。注意大熊座中构成北斗七星的七颗恒星，并从北斗最外侧的两颗恒星（有时称为“指极星”）连成一条直线指向北极星 Polaris。} \label{fig_Ursa_3}
\end{figure}
\subsubsection{恒星}  
德国制图学家约翰·拜耳使用希腊字母 $\alpha$ 至 $\theta$ 来标注该星座中最显著的恒星，而他的同乡约翰·埃勒特·博德随后又补充使用了 $\iota$ 至 $\phi$。目前仅有 $\lambda$ 与 $\pi$ 仍在使用，这很可能是由于它们靠近北天极的缘故 $^{[16]}$。在该星座的边界之内，共有 $39$ 颗视星等亮于或等于 $6.5$ 等的恒星 $^{[22][c]}$。  

按照约翰·拜耳的排序，主要七颗恒星的传统名称为：  
\begin{enumerate}
\item Polaris  
\item Kochab  
\item Pherkad  
\item Yildun  
\item Epsilon Ursae Minoris 无传统名称。  
\item Zeta Ursae Minoris 无传统名称。  
\item Eta Ursae Minoris 无传统名称。  
\end{enumerate}
作为小熊尾部的标志 $^{[16]}$，北极星 Polaris（即小熊座 $\alpha$ 星）是该星座中最亮的恒星，其视星等在 $1.97$ 与 $2.00$ 之间变化，周期为 $3.97$ 天 $^{[24]}$。它距离地球约 $432$ 光年 $^{[25]}$，是一颗黄白色超巨星，光谱型在 F7Ib 与 F8Ib 之间变化 $^{[24]}$，其质量约为太阳的 $6$ 倍，光度约为太阳的 $2500$ 倍，半径约为太阳的 $45$ 倍。北极星是从地球可见的最亮的造父变星。它是一个三星系统，其中的超巨星主星拥有两颗黄白色主序星伴星，轨道半长轴分别约为 $17$ 和 $2400$ 天文单位（AU），其公转周期分别为 $29.6$ 年与 $42000$ 年 $^{[26]}$。  

传统名称为 Kochab 的小熊座 $\beta$ 星，其视星等为 $2.08$，略暗于北极星 $^{[27]}$。它距离地球约 $131$ 光年 $^{[28][d]}$，是一颗橙色巨星——即核心氢已耗尽、离开主序阶段的演化恒星——光谱型为 K4III $^{[27]}$。Kochab 在约 $4.6$ 天的周期内呈现轻微的光变，其质量通过这些振荡测量估计约为太阳的 $1.3$ 倍 $^{[29]}$。Kochab 的光度约为太阳的 $450$ 倍，直径约为太阳的 $42$ 倍，表面温度约为 $4130\,\mathrm{K}$ $^{[30]}$。其年龄估计约为 $29.5$ 亿年，误差约 $\pm 10$ 亿年，并已被宣布拥有一颗行星伴星，其质量约为木星的 $6.1$ 倍，轨道周期为 $522$ 天 $^{[31]}$。
\begin{figure}[ht]
\centering
\includegraphics[width=6cm]{./figures/58bc21206ce95031.png}
\caption{小熊座与大熊座相对于北极星 Polaris 的位置关系} \label{fig_Ursa_4}
\end{figure}
